\section{Approccio generale}
Come prima cosa si è deciso di affrontare il problema in maniera del rilevamento degli accenti a partire da un testo in lingua italiana.
Per riuscire ad arrivare a un modello capace di riuscire nel compito dell'induviazione degli accenti è necessario importaqre dei dataset importati per il task. Fatto ciò tali dataset saranno usati per creare per appunto il modello che data una generica frase in italiano possa individuare tutti gli accenti con correttezza
\subsection{Q2Stress}
Q2Stress è un dataset di frasi in lingua italiana, ognuna delle quali è annotata con informazioni sugli accenti presenti. Questo dataset è stato utilizzato per addestrare il modello di riconoscimento degli accenti, fornendo esempi concreti di come gli accenti si manifestano nel linguaggio naturale.
\subsection{phonItaliaR}
phonItaliaR è un dataset di frasi in lingua italiana, ognuna delle quali è annotata con informazioni sulle pronunce e sugli accenti presenti. Questo dataset è stato utilizzato per addestrare il modello di riconoscimento degli accenti, fornendo esempi concreti di come gli accenti si manifestano nel linguaggio naturale.
\section{Definizione del modello}
\begin{definition}
Un modello matematico \`e una rappresentazione semplificata di un fenomeno reale, costruita per descrivere relazioni quantificabili tra variabili.
\end{definition}

Per esempio, una funzione di costo pu\`o essere definita come
\begin{equation}
\mathcal{L}(\theta) = \frac{1}{n} \sum_{i=1}^{n} \ell(y_i, f_\theta(x_i)) + \lambda \|\theta\|_2^2.
\end{equation}

\section{Ottimizzazione}
L'ottimizzazione di questa espressione richiede tecniche numeriche, come gradient descent o metodi quasi-Newton, a seconda della complessit\`a del problema e della dimensione dello spazio dei parametri.

\begin{remark}
La scelta della funzione di perdita influenza in modo diretto la robustezza del modello rispetto a rumore, outlier e distribuzioni non gaussiane.
\end{remark}
